% Dokument / Seiten
\documentclass[12pt,a4paper]{scrartcl}
\usepackage[utf8]{inputenc}				% Zeichencodierung
\usepackage[ngerman]{babel}				% Sprachpaket 
\usepackage[T1]{fontenc}				% Umlaute bei \hyphenation verwenden können (Silbentrennung)
\usepackage[onehalfspacing]{setspace}	% Setzt das Dokument auf 1,5x Zeilenabstand

\usepackage{							% standard Krempel
	amsmath,
	amsfonts,
	amsthm,
	amssymb}

% Gleitumgebungen / Grafiken
\usepackage{float}						% notwendig für Gleitubgebungen z.B. Figure
\usepackage{subfig}						% erlaubt das Verwenden von Subfigures
\usepackage{graphicx}					% erlaubt das Einbinden von Grafiken
\usepackage{pdfpages}					% ermöglicht das Einbinden von PDFs
\usepackage{epstopdf}					% einbinden von eps Bilddateien
\usepackage{longtable}					% lange Tabellen verarbeiten
\usepackage{caption}					% ermöglicht die Beschriftung von Gleitobjekten
\usepackage{nameref}					% erlaubt Namensreferenzen

% Diagramme
\usepackage{tikz}						% Zeichentool
\usepackage{pgfplots, pgfplotstable}	% Diagramme zeichnen
\pgfplotsset{compat = newest}			% Einstellung für Diagramme
\usepackage{siunitx}                    % SI-Einheiten
\usepgfplotslibrary{units}				% Einheiten in Diagrammen plotten

% Verzeichnisse
\usepackage[square,numbers]{natbib}		% Literatur zitieren
\bibliographystyle{dinat}				% Literatur- und Zitatstyle

\usepackage{makeidx}					% anlegen eines automatischen Inhaltsverzeichnis
\usepackage{tocstyle}					% Punkte im Inhaltsverzeichnis
\usetocstyle{allwithdot} 

% Formatierungen
\setlength{\parindent}{0in} 			% Absatz, erste Zeile wird global NICHT eingerückt
\usepackage{parskip}					% setzt einen Abstand nach einem \par Absatz

\usepackage[headsepline]{scrpage2}		% Kopfzeile

\usepackage[							% Verlinkung aller \ref,url,etc.
colorlinks=true,
urlcolor=blue,
citecolor=blue,
linkcolor=black
]{hyperref}								

\usepackage{wasysym}					% zusätzliche Symbole verwenden (Checkbox)

% Silbentrennung
\expandafter\def\expandafter\UrlBreaks\expandafter{\UrlBreaks\do\a		% Erlaubt einen Umbruch in einer URL 
	\do\b\do\c\do\d\do\e\do\f\do\g\do\h\do\i\do\j\do\k\do\l\do\m\do\n	% an allen angegebenen Stellen
	\do\o\do\p\do\q\do\r\do\s\do\t\do\u\do\v\do\w\do\x\do\y\do\z\do\A	
	\do\B\do\C\do\D\do\E\do\F\do\G\do\H\do\I\do\J\do\K\do\L\do\M\do\N	
	\do\O\do\P\do\Q\do\R\do\S\do\T\do\U\do\V\do\W\do\X\do\Y\do\Z\do\1
	\do\2\do\3\do\4\do\5\do\6\do\7\do\8\do\9\do\0\do\&} 	 

%\sloppy								% Nachlässige Silbentrennung, Blocksatzoptimiert
\hyphenation{OpenOCD}					% Worte, die NICHT oden nur auf bestimmte Weise getrennt werden sollen

% Schusterjungen und Hurenkinder
\usepackage[all]{nowidow}
%\clubpenalty = 10000 % schliesst Schusterjungen aus
%\widowpenalty = 10000 
%\displaywidowpenalty = 10000% schliesst Hurenkinder aus 


% Definition von Funktionen

% \paragraph mit Zeilenumbruch versehen ---------------------------
\makeatletter
\renewcommand\paragraph{\@startsection{paragraph}{4}{\z@}%
	{-3.25ex\@plus -1ex \@minus -.2ex}%
	{1.5ex \@plus .2ex}%
	{\normalfont\normalsize\bfseries}}
\makeatother
%------------------------------------------------------------------



% Beginn des Dokuments --------------------------------------------
\begin{document}
	
	\ohead{\pagemark}
	\ihead{\sectionmark}
	\cfoot{}
	\pagestyle{scrheadings}
	\renewcommand{\sectionmark}[1]{\markright{#1}{}}
	
	% Titelseite ------------------------------------------------------------------
	% Titelseite ------------------------------------------------------------------
	\begin{titlepage}
		
		\begin{center}
			\includegraphics[height=4cm]{Bilder/HTWK_Logo_CMYK}
			
			\vspace*{3cm}
			
			{\LARGE Ein wirklich gutes Thema\\ auf zwei Zeilen}\\
			\vspace{1cm}
			{\Large Belegarbeit}\\
		\end{center}
		
		\vspace{8cm}
		
		\begin{tabbing}
			\hspace{4cm}\=\kill
			Vorgelegt von:		\> Mario Hesse \\
			Eingereicht bei:	\> \dots \\ 
			Im Studienfach:		\> \dots \\
			Ort,Datum:			\> Leipzig, den \today
		\end{tabbing} 
		
	\end{titlepage}
	
	% Sperrvermerk ------------------------------------------------------------------------
	
	%\thispagestyle{empty} % no page number
	%\section*{Sperrvermerk}
	%
	%\medskip
	%\noindent
	%\emph{
	%Die vorliegende Bachelorarbeit beinhaltet interne vertrauliche Informationen des Laboratory for Biosignal Processing (LaBP). Die Weitergabe des Inhaltes der Arbeit und eventuell beiliegender
	%Zeichnungen und Daten im Gesamten oder in Teilen ist grundsätzlich untersagt. Es dürfen keinerlei Kopien oder Abschriften -- auch in digitaler Form -- gefertigt werden. Ausnahmen bedürfen der vorherigen schriftlichen Genehmigung des Laboratory for Biosignal Processing (LaBP).\\[1cm]
	%Laboratory for Biosignal Processing\\
	%Dr.-Ing. Gerold Bausch\\
	%Eilenburger Straße 13\\
	%04317 Leipzig\\
	%Telefon: +49 341 3076 3103 \\
	%E-Mail: gerold.bausch@htwk-leipzig.de\\[2cm]
	%Leipzig, den 3.Oktober 2015\\[2,5cm]
	%\begin{tabbing}
	%\hspace*{3cm}   \= \hspace*{6cm} \= \kill
	% \> Gerold Bausch \> Mario Hesse \\	
	%\end{tabbing}
	%}
	
	
	
	
	% Vorwort, Danksagung ------------------------------------------------------------------------
	\clearpage
	\phantomsection
	\addcontentsline{toc}{section}{Vorwort}
	\section*{Vorwort}
	
	Persönliche Verbindung zum Thema. Grund für diesen Artikel. 
	
	
	% Danksagung ------------------------------------------------------------------------------------
	\clearpage
	\phantomsection
	\addcontentsline{toc}{section}{Danksagung}
	\section*{Danksagung}
	
	Bei größeren Arbeiten bedankt man sich bei Menschen und Institutionen, die Maßgeblich an der Entstehung teil hatten. (beim Bedanken Rangfolge beachten, wem gebührt der größte Dank) \\
	
	
	% Kurzfassung, Abstract ------------------------------------------------------------------------
	\clearpage
	\phantomsection
	\addcontentsline{toc}{section}{Kurzfassung, Abstract}
	\section*{Kurzfassung}
	
	Hier wird der Inhalt der Arbeit in ein paar Zeilen kurz dargestellt. Man soll auf Grund dieser Angabe entscheiden können, ob es lohnt die Arbeit zu lesen oder nicht.
	
	\medskip
	\noindent
	Schlagwörter: Überblick, Zusammenfassung, Einblick \\
	
	
	\section*{Abstract}
	
	This is the same than in "Kurzfassung" just in english language. This part is used to internationalize an article. 
	
	\medskip
	\noindent
	Keywords: summary, content, overview \\
	
	
	% Inhaltsverzeichnis ------------------------------------------------------------------------
	\clearpage
	\phantomsection
	\addcontentsline{toc}{section}{Inhaltsverzeichnis}
	\tableofcontents
	
	
	% Abbildungsverzeichnis ------------------------------------------------------------------------
	\clearpage
	\phantomsection
	\addcontentsline{toc}{section}{Abbildungsverzeichnis}
	\listoffigures
	
	
	% Tabellenverzeichnis ------------------------------------------------------------------------
	\clearpage
	\phantomsection
	\addcontentsline{toc}{section}{Tabellenverzeichnis}
	\listoftables
	
	
	
	\vfill
	% Abkürzungsverzeichnis ------------------------------------------------------------------------
	\phantomsection
	\addcontentsline{toc}{section}{Abkürzungsverzeichnis}
	\section*{Abkürzungsverzeichnis}
	
	
	\begin{tabbing}
		\hspace{2cm}\=\kill
		HTWK 	\> Hochschule für Technik, Wirtschaft und Kultur\\
		k.A. 	\> keine Angabe \\
	\end{tabbing} 
	
	\vfill
	
	% Einheitenverzeichnis --------------------------------------------------------------------------
	\clearpage
	\phantomsection
	\addcontentsline{toc}{section}{Einheitenverzeichnis}
	\section*{Einheitenverzeichnis}
	
	\begin{tabbing}
		\hspace{2cm}\=\kill	
		
		$V$			\> Volt \\
		$A$			\> Ampere \\
		$\Omega$	\> Ohm\\
		
	\end{tabbing} 
	
	
	% Formelverzeichnis ------------------------------------------------------------------------
	\clearpage
	\phantomsection
	\addcontentsline{toc}{section}{Formelverzeichnis}
	\section*{Formelverzeichnis}
	
	\begin{tabbing}
		\hspace{2cm}\=\kill
		
		$U$					\> Spannung (allgemein)\\
		$R$					\> Widerstand (allgemein)\\
		$I$					\> Strom (allgemein)\\
		
	\end{tabbing} 
	
	
	% Glossar ---------------------------------------------------------------------------------
	\clearpage
	\phantomsection
	\addcontentsline{toc}{section}{Glossar}
	\section*{Glossar}
	
	\begin{longtable}{p{4cm}p{11cm}}
		
		
		Glossar			& \dots ist eine Tabelle in der Fachwörter für die forliegenede Arbeit definiert sind.\\
	\end{longtable}
	
	\addtocounter{table}{-1}	% Die Tabelle "Glossar wird nicht mit ins Tabellenverzeichnis aufgenommen
	
	% Dokumententext ------------------------------------------------------------------------
	
	% Literatur  -----------------------------------------------------------------
	\clearpage
	\phantomsection
	\addcontentsline{toc}{section}{Literaturverzeichnis}
	\sectionmark{Literaturverzeichnis}
	\bibliography{Literatur} % nutzt Angaben in der Datei "Literatur.bib"
	
	
	% Eidesstattliche Versicherung ------------------------------------------------------------------------
	\clearpage
	\thispagestyle{empty}
	\phantomsection
	\addcontentsline{toc}{section}{Eidesstattliche Versicherung}
	\section*{Eidesstattliche Erklärung}
	
	\smallskip
	Ich erkläre hiermit, dass ich diese Bachelorarbeit selbstständig ohne Hilfe Dritter und ohne Benutzung anderer als der angegebenen Quellen und Hilfsmittel verfasst habe. Alle den benutzten Quellen wörtlich oder sinngemäß entnommenen Stellen sind als solche einzeln kenntlich gemacht.
	
	\smallskip
	\noindent
	Diese Arbeit ist bislang keiner anderen Prüfungsbehörde vorgelegt und auch nicht veröffentlicht worden.
	
	\smallskip
	\noindent
	Ich bin mir bewusst, dass eine falsche Erklärung rechtliche Folgen haben wird.
	
	\bigskip
	\noindent
	Leipzig, den 
	
	\pagestyle{empty}
	\newpage % Leerseite für Notizen
	~
	
	% ***********************************************
\end{document}
% ***********************************************

